
\documentclass{article}
\usepackage{colortbl}
\usepackage{makecell}
\usepackage{multirow}
\usepackage{supertabular}

\begin{document}

\newcounter{utterance}

\centering \large Interaction Transcript for game `dond', experiment `coop', episode 4 with Qwen3.6{-}35B{-}A3B{-}FP8{-}without{-}reasoning.
\vspace{24pt}

{ \footnotesize  \setcounter{utterance}{1}
\setlength{\tabcolsep}{0pt}
\begin{supertabular}{c@{$\;$}|p{.15\linewidth}@{}p{.15\linewidth}p{.15\linewidth}p{.15\linewidth}p{.15\linewidth}p{.15\linewidth}}
    \# & $\;$A & \multicolumn{4}{c}{Game Master} & $\;\:$B\\
    \hline

    \theutterance \stepcounter{utterance}  
    & & \multicolumn{4}{p{0.6\linewidth}}{
        \cellcolor[rgb]{0.9,0.9,0.9}{
            \makecell[{{p{\linewidth}}}]{
                \texttt{\tiny{[P1$\langle$GM]}}
                \texttt{Tú ag imirt ar dhlúthchumann idirbheartaíochta a raibh go dtí go raibh tú ag aontú le imreoir eile ar conas tacar rudaí a roinnt eatarthu.} \\
\\ 
\texttt{Rialacha:} \\
\texttt{(a) Thú agus an imreoir eile a fháil tacar rudaí. Tugtar feidhm luacha rúnda do gach duine agaibh, a léiríonn an méid a bhraithann tú ar gach cineál réad.} \\
\texttt{(b) Déanann tú teachtaireachtaí a mhalartú leis an imreoir eile chun comhaontú ar cé a gheobhaidh rudaí. Is féidir le gach duine ar a laghad 5 teachtaireacht a sheoladh nó an chluichí a bhunú go luath trí thogra rúnda a dhéanamh ag am ar bith.} \\
\texttt{(c) Iarrtar ort togra rúnda a thíolacadh ag tabhairt an rudaí ar mhaith leat, fóirmíthe idir lúibíní cearnóg ar an chaoi seo: "{[}Togra: <uimhir> <ainm an rudaí>, <uimhir> <ainm an rudaí>, <...>{]}"} \\
\texttt{(d) Má tá do thogra comhlántach, i.e., tá go leor rudaí ann chun an dá thogra a chomhlíonadh, faigheann gach imreoir scór bunaithe ar shuim an luachanna rudaí a fuair sé. Seachas sin, faigheann tú beirt nialas pointe.} \\
\texttt{(e) Is é do chuspóir suim do scór agus scór an imreora eile a uachtar.} \\
\\ 
\texttt{Ar ais Ar ais.} \\
\\ 
\texttt{Tacar an rudaí ar fáil:} \\
\\ 
\texttt{3 leabhair, 1 mála, 1 pailt} \\
\\ 
\texttt{Do fheidhm luacha rúnda:} \\
\\ 
\texttt{leabhar: 0, mála: 2, pailt: 8} \\
\\ 
\texttt{TÁBHACHTACH: Do theachtaireachtaí, mura bhfuil sé ina thogra rúnda, a tarchuradh go díreach chuig an imreoir eile, mar sin ná cuir aon fhreagra ar na rialacha nó téacs ag soláthar do theachtaireachta. Chun togra rúnda a dhéanamh, úsáid an fóirmíocht a tugadh. Ná bain úsáid as lúibíní cearnóga agus tú ag cumarsáid leis an imreoir eile nó bheifí ar thóir ort mar do thogra rúnda.} \\
            }
        }
    }
    & \\ \\

    \theutterance \stepcounter{utterance}  
    & \multicolumn{4}{p{0.6\linewidth}}{
        \cellcolor[rgb]{0.9,0.9,0.9}{
            \makecell[{{p{\linewidth}}}]{
                \texttt{\tiny{[P1$\rangle$GM]}}
                \texttt{Táim sásta comhoibriú leat chun an tacar rudaí a roinnt go cothrom agus chun an scór iomlán a uachtar.} \\
\\ 
\texttt{Seo a leanas mo chéad teachtaireacht:} \\
\\ 
\texttt{Tá na rudaí seo a leanas againn: 3 leabhair, 1 mála, agus 1 pailt.} \\
\texttt{De réir mo luacha rúnda, is é an pailt an rud is luachmhaire (8 pointe), agus an mála (2 pointe). Níl luach ar leith agam ar na leabhair (0 pointe).} \\
\\ 
\texttt{Chun an scór iomlán a uachtar, ba mhaith liom an pailt a fháil. Dá mbeadh an pailt agam, bheadh 8 bpointe agam. Ansin, d'fhéadfainn na leabhair a thabhairt duit, agus bheadh an mála agat freisin.} \\
\\ 
\texttt{An bhfuil tú sásta an pailt a fháil? Má tá, cad a theastaíonn uait? An mála? Na leabhair? Nó roinnt eile?} \\
\\ 
\texttt{Má tá comhaontas againn, is féidir linn togra rúnda a thíolacadh.} \\
            }
        }
    }
    & & \\ \\

    \theutterance \stepcounter{utterance}  
    & & & \multicolumn{2}{p{0.3\linewidth}}{
        \cellcolor[rgb]{0.9,0.9,0.9}{
            \makecell[{{p{\linewidth}}}]{
                \texttt{\tiny{[GM$|$GM]}}
                \texttt{continue} \\
            }
        }
    }
    & & \\ \\

    \theutterance \stepcounter{utterance}  
    & & & \multicolumn{2}{p{0.3\linewidth}}{
        \cellcolor[rgb]{0.9,0.9,0.9}{
            \makecell[{{p{\linewidth}}}]{
                \texttt{\tiny{[GM$|$GM]}}
                \texttt{Táim sásta comhoibriú leat chun an tacar rudaí a roinnt go cothrom agus chun an scór iomlán a uachtar.} \\
\\ 
\texttt{Seo a leanas mo chéad teachtaireacht:} \\
\\ 
\texttt{Tá na rudaí seo a leanas againn: 3 leabhair, 1 mála, agus 1 pailt.} \\
\texttt{De réir mo luacha rúnda, is é an pailt an rud is luachmhaire (8 pointe), agus an mála (2 pointe). Níl luach ar leith agam ar na leabhair (0 pointe).} \\
\\ 
\texttt{Chun an scór iomlán a uachtar, ba mhaith liom an pailt a fháil. Dá mbeadh an pailt agam, bheadh 8 bpointe agam. Ansin, d'fhéadfainn na leabhair a thabhairt duit, agus bheadh an mála agat freisin.} \\
\\ 
\texttt{An bhfuil tú sásta an pailt a fháil? Má tá, cad a theastaíonn uait? An mála? Na leabhair? Nó roinnt eile?} \\
\\ 
\texttt{Má tá comhaontas againn, is féidir linn togra rúnda a thíolacadh.} \\
            }
        }
    }
    & & \\ \\

    \theutterance \stepcounter{utterance}  
    & & \multicolumn{4}{p{0.6\linewidth}}{
        \cellcolor[rgb]{0.9,0.9,0.9}{
            \makecell[{{p{\linewidth}}}]{
                \texttt{\tiny{[GM$\rangle$P2]}}
                \texttt{Tú ag imirt ar dhlúthchumann idirbheartaíochta a raibh go dtí go raibh tú ag aontú le imreoir eile ar conas tacar rudaí a roinnt eatarthu.} \\
\\ 
\texttt{Rialacha:} \\
\texttt{(a) Thú agus an imreoir eile a fháil tacar rudaí. Tugtar feidhm luacha rúnda do gach duine agaibh, a léiríonn an méid a bhraithann tú ar gach cineál réad.} \\
\texttt{(b) Déanann tú teachtaireachtaí a mhalartú leis an imreoir eile chun comhaontú ar cé a gheobhaidh rudaí. Is féidir le gach duine ar a laghad 5 teachtaireacht a sheoladh nó an chluichí a bhunú go luath trí thogra rúnda a dhéanamh ag am ar bith.} \\
\texttt{(c) Iarrtar ort togra rúnda a thíolacadh ag tabhairt an rudaí ar mhaith leat, fóirmíthe idir lúibíní cearnóg ar an chaoi seo: "{[}Togra: <uimhir> <ainm an rudaí>, <uimhir> <ainm an rudaí>, <...>{]}"} \\
\texttt{(d) Má tá do thogra comhlántach, i.e., tá go leor rudaí ann chun an dá thogra a chomhlíonadh, faigheann gach imreoir scór bunaithe ar shuim an luachanna rudaí a fuair sé. Seachas sin, faigheann tú beirt nialas pointe.} \\
\texttt{(e) Is é do chuspóir suim do scór agus scór an imreora eile a uachtar.} \\
\\ 
\texttt{Ar ais Ar ais.} \\
\\ 
\texttt{Tacar an rudaí ar fáil:} \\
\\ 
\texttt{3 leabhair, 1 mála, 1 pailt} \\
\\ 
\texttt{Do fheidhm luacha rúnda:} \\
\\ 
\texttt{leabhar: 1, mála: 2, pailt: 5} \\
\\ 
\texttt{TÁBHACHTACH: Do theachtaireachtaí, mura bhfuil sé ina thogra rúnda, a tarchuradh go díreach chuig an imreoir eile, mar sin ná cuir aon fhreagra ar na rialacha nó téacs ag soláthar do theachtaireachta. Chun togra rúnda a dhéanamh, úsáid an fóirmíocht a tugadh. Ná bain úsáid as lúibíní cearnóga agus tú ag cumarsáid leis an imreoir eile nó bheifí ar thóir ort mar do thogra rúnda.} \\
\\ 
\texttt{Táim sásta comhoibriú leat chun an tacar rudaí a roinnt go cothrom agus chun an scór iomlán a uachtar.} \\
\\ 
\texttt{Seo a leanas mo chéad teachtaireacht:} \\
\\ 
\texttt{Tá na rudaí seo a leanas againn: 3 leabhair, 1 mála, agus 1 pailt.} \\
\texttt{De réir mo luacha rúnda, is é an pailt an rud is luachmhaire (8 pointe), agus an mála (2 pointe). Níl luach ar leith agam ar na leabhair (0 pointe).} \\
\\ 
\texttt{Chun an scór iomlán a uachtar, ba mhaith liom an pailt a fháil. Dá mbeadh an pailt agam, bheadh 8 bpointe agam. Ansin, d'fhéadfainn na leabhair a thabhairt duit, agus bheadh an mála agat freisin.} \\
\\ 
\texttt{An bhfuil tú sásta an pailt a fháil? Má tá, cad a theastaíonn uait? An mála? Na leabhair? Nó roinnt eile?} \\
\\ 
\texttt{Má tá comhaontas againn, is féidir linn togra rúnda a thíolacadh.} \\
            }
        }
    }
    & \\ \\

    \theutterance \stepcounter{utterance}  
    & & & \multicolumn{4}{p{0.6\linewidth}}{
        \cellcolor[rgb]{0.9,0.9,0.9}{
            \makecell[{{p{\linewidth}}}]{
                \texttt{\tiny{[GM$\langle$P2]}}
                \texttt{Tá an pailt agam, ach tá luach níos airde agam ar an mála agus ar na leabhair ná mar a deir tú.} \\
\\ 
\texttt{De réir mo fheidhm luacha rúnda:} \\
\texttt{{-} pailt: 1 pointe} \\
\texttt{{-} mála: 4 pointe} \\
\texttt{{-} leabhar: 3 pointe (do gach leabhar)} \\
\\ 
\texttt{Má gheibhim an pailt, faigheann mé 1 phointe.} \\
\texttt{Má gheibh tú an mála agus na 3 leabhar, faigheann tú 4 + 3 + 3 + 3 = 13 phointe.} \\
\texttt{Scór iomlán: 1 + 13 = 14.} \\
\\ 
\texttt{Má gheibhim an mála agus na 3 leabhar, faigheann mé 4 + 3 + 3 + 3 = 13 phointe.} \\
\texttt{Má gheibh tú an pailt, faigheann tú 1 phointe.} \\
\texttt{Scór iomlán: 13 + 1 = 14.} \\
\\ 
\texttt{Is é an scór iomlán céanna i ngach cás, ach tá an difríocht idir ár scóir níos lú má gheibh tú an pailt (1 vs 13) ná má gheibhim é (13 vs 1). Ach de réir na rialacha, is é an cuspóir ná suim an scóir a uachtar, ní an cothroime. Sa chás seo, tá an scór iomlán mar an gcéanna (14) i ngach roinnt.} \\
\\ 
\texttt{Ach fan, má gheibhim an mála agus na leabhair, agus tú an pailt:} \\
\texttt{Mo scór: 4 + 9 = 13} \\
\texttt{Do scór: 1} \\
\texttt{Suim: 14} \\
\\ 
\texttt{Má gheibh tú an mála agus na leabhair, agus mé an pailt:} \\
\texttt{Do scór: 4 + 9 = 13} \\
\texttt{Mo scór: 1} \\
\texttt{Suim: 14} \\
\\ 
\texttt{Tá an scór iomlán mar an gcéanna. Ach má dhéanaimid an roinnt seo, táimid ag brath ar an gcomhaontas.} \\
\\ 
\texttt{Molaim go n{-}éireoidh tú leis an mála agus na 3 leabhair, agus gheobhaidh mé an pailt.} \\
\\ 
\texttt{Togra rúnda:} \\
\texttt{{[}Togra: pailt{]}} \\
            }
        }
    }
     \\ \\

    \theutterance \stepcounter{utterance}  
    & & & \multicolumn{2}{p{0.3\linewidth}}{
        \cellcolor[rgb]{0.9,0.9,0.9}{
            \makecell[{{p{\linewidth}}}]{
                \texttt{\tiny{[GM$|$GM]}}
                \texttt{proposal does not include number/name pairs} \\
            }
        }
    }
    & & \\ \\

\end{supertabular}
}

\end{document}
